\subsection{The first new degree: vanishing of \texorpdfstring{$K_{12}$}{K12}}
\label{sec:k12-hecke-vanishing}

\begin{theorem}[K17-K12-VANISH]\label{thm:k12-vanish}
Under the stipulated arithmetic theorems of Alp\"oge's source,
\[
                         K_{12}(\mathbf Z)=0.
\]
Consequently, for every odd prime $\ell$,
\[
 A_\ell^{[-6]}=0,
\]
and the full-domain Soul\'e character $\kappa_7$ has image
$\mathbf Z_\ell$ in the normalization of VK-SOULE-INDEX.
\end{theorem}

\begin{proof}
The 2-primary subgroup vanishes by the Rognes--Weibel comparison already
cited in (K4). For every odd $\ell<2^{31}$, the published Vandiver
verification of Hart--Harvey--Ong, together with (K2), gives
$K_{12}(\mathbf Z)\{\ell\}=0$. This verification is cited, not rerun.

Now let $\ell>2^{31}$ be prime. Take the auxiliary field
$E=\mathbf Q(\sqrt{-3})$, the twist $e=6$, the modulus $N=2$ and the
trivial Dirichlet character. All the source datum conditions in K17-S17-Q
hold: $E$ has class number one, $\ell\ne2,3$, the required parity is even,
and $(\ell-1)\nmid12$. The exact specialization is also the one in
\texttt{SmallClassNumbers/ClassNumberThree.lean}.

The prime seven splits in $E$. Choose arithmetic Frobenius $\sigma$
at the factor determined by $3+\omega$, and set $s=\sigma^{-1}$.
The comparison trace in Lemma~\ref{lem:step17-unit} is
\[
                t=7^{-6}+1+7^6.
\]
For every $\mathrm{Adm3U\_s}(\iota,6,2)$ representation,
Lemma~\ref{lem:step17-u3-level}, Theorem~\ref{thm:u3-coverage}, and
Theorem~\ref{thm:u3-hecke} give
\[
 \iota\operatorname{tr}\rho(s)=\frac{157899}{7^6}.
\]
For every $\mathrm{Adm2A}(\iota,6,6)$ representation,
Lemma~\ref{lem:gl2-conductor} places its trace in
$1+a_7(f)/7^6$ for an eigenform in the complete space of
Proposition~\ref{prop:gl2-weight13}.

Put $E_7=7^{12}+7^6+1$. The two exact certificates give the nonzero
integer
\begin{align*}
 N_7
 &=(E_7-157899)P_7(1+7^{12})\\
 &=36703455074769833036351681539319447519232\\
 &=2^{15}3^{19}\,7\,13^5\,37\,43^3\,47^3\,61\,271^3.
\end{align*}
Every prime factor is at most $271$, so $\ell\nmid7N_7$.
The irreducible-branch trace difference is therefore an $\ell$-adic
unit. In the rank-two branch the eigenvalue $a_7(f)$ is an algebraic
integer and a root of the monic integral polynomial $P_7$.
Since $P_7(1+7^{12})$ is a unit at every prime above $\ell$, every factor
$1+7^{12}-a_7(f)$ is a unit there. Dividing by $7^6$ preserves this
property. Thus every admissible trace difference is an $\ell$-adic unit.

Lemma~\ref{lem:step17-unit} forces the source Selmer length to be zero.
The rational-field identification in K17-S17-Q and the exact strict
duality (K5) give $K_{12}(\mathbf Z)\{\ell\}=0$. Every prime is now
covered. Finiteness of $K_{12}(\mathbf Z)$ proves its vanishing.

The class-group assertion follows from (K2) with $k=3$.
For $m=7$, the full Soul\'e image index is
$|K_{2m-2}(\mathbf Z)\{\ell\}|=1$ by VK-SOULE-INDEX, proving the
last assertion.
\end{proof}

\begin{corollary}[K17-K12-PERIOD]\label{cor:k12-period}
For every odd prime $\ell$ and every positive $k$ satisfying
$k\equiv3\pmod{(\ell-1)/2}$, one has
$K_{4k}(\mathbf Z)\{\ell\}=0$.
Together with the background cases, the residue classes
$0,1,2,3$ are therefore excluded for all odd primes.
\end{corollary}
\begin{proof}
The character $-2k$ modulo $\ell-1$ equals $-6$ in the stated residue
class. Apply (K2) and Theorem~\ref{thm:k12-vanish}; the residues
$0,1,2$ were treated after (K4).
\end{proof}

\paragraph{Scope of the new degree.}
The proof uses the stronger intermediate Selmer inequality with the
actual source predicates. It neither infers vanishing from an
unspecified numerical bound nor replaces the old oversized trace
product by a divisor. It closes the numerical $K_{12}$ question without
requiring any of the remaining Steinberg boundary matrices.
The individual component $A_\ell^{[-6]}$ does not imply all of Vandiver
at an arbitrary prime. Likewise, surjectivity of the full-domain
$\kappa_7$ does not remove the factorial restriction on its weight
subgroup or the integral even-weight defects at irregular primes.
